\scanpage{32}
\mspage{1}
I) Arrangements de $3$ pseudo-droites

1) Cube combinatoire (de dim 3). La donnée équivaut à celle d'une famille $(\Phi_i)_{i\in I}$ de trois ens à deux éléments ($\operatorname{card} I=3$, $\operatorname{card}\Phi_i=2\quad\forall i\in I$)
\diagram{(1)\qquad
\begin{tikzcd}[ampersand replacement=\&,column sep=large,row sep=small]
\Phi=\coprod\limits_{i\in I}\Phi_i \arrow[d] \& \text{card }6\\
I \& \text{card }3
\end{tikzcd}}

Les éléments de $\Phi$ s'identifient aux faces du cube, ceux de $I$ aux paires de faces opposées. Les facettes du cube s'identifient aux \uline{sections partielles} de $\Phi$ sur \inserted{des parties $\ne\varnothing$ de} $I$, avec la relation d'incidence correspondant à la relation \uline{opposée} à celle du prolongement de sections partielles, i.e. de l'inclusion des parties correspondantes de $\Phi$ (parties $A\subset\Phi$ non vides, telles que l'appl. induite $A\to I$ soit injective).
\[
\begin{aligned}
 C_2&\simeq\Phi\simeq\mathfrak P_1(\Phi)\\
 C_1&\simeq\{A\in\mathfrak P_2(\Phi)\mid A\to I\text{ injectif}\}\\
 C_0&\simeq\{A\in\mathfrak P_3(\Phi)\mid A\to I\text{ injectif}\ (\text{donc bij.})\}
\end{aligned}\tag*{(2)}
\]
On a \unsure{une|un} automorphisme \inserted{(involutif)} canonique $\underline a$ de (1) induisant l'identité dans $I$ et la transposition dans chacun des $\Phi_i$, qui induit un autom. \inserted{(central)}, encore noté $\underline a$, du cube (2), lequel est l'antipodisme.

Si on prend le polyèdre opposé au cube, on trouve l'octaèdre donné par
\[
\left\{\begin{aligned}
 O_0&=C_2\simeq\Phi\simeq\mathfrak P_1(\Phi)\\
 O_1&=C_1\simeq\{A\in\mathfrak P_2(\Phi)\mid A\hookrightarrow I\}\\
 O_2&=C_0\simeq\{A\in\mathfrak P_3(\Phi)\mid A\unsure{\hookrightarrow|\lhook\joinrel\relbar} I\}
\end{aligned}\right.\tag*{(3)}
\]
avec cette fois-ci la relation \inserted{d'incidence qui est} de prolongement des sections
% archive sheet number: 31
